% !TEX root = ../main.tex
% ==============================================================================
% SECCIÓN 4: PRESUPUESTO DETALLADO Y VIABILIDAD FINANCIERA
% Responsable: Joaquín (Líder Técnico / Complejidad Alta)
% Proyecto: Campus Zero - Control Automatizado de Iluminación y Climatización
% En coordinación directa con: Javier (Justificaciones y Cotizaciones)
% ==============================================================================

\section{Presupuesto Detallado y Viabilidad Financiera}
\label{sec:presupuesto}

\subsection{Estrategia Financiera y Coherencia con Fondos Institucionales USM}

La viabilidad material de la iniciativa \textit{Campus Zero: Control Automatizado de Iluminación y Climatización de Salas de Estudio} exige un diseño presupuestario gobernado por la rigurosidad técnica, la optimización extrema del gasto y el estricto apego normativo de la Universidad Técnica Federico Santa María (UTFSM). El proyecto se postula formalmente a la línea de \textbf{Sustentabilidad y Medio Ambiente} de los Fondos Concursables Institucionales de la USM (Dirección General de Infraestructura, DGI, y Dirección de Vinculación con el Medio), concebidos para impulsar soluciones tecnológicas desarrolladas por la comunidad académica que reduzcan la huella de carbono y mejoren el uso de recursos energéticos en los campus universitarios.

Las bases concursables establecen un techo presupuestario máximo de \textbf{\$1.250.000~CLP}. En estricta sujeción a este marco, el equipo formuló una asignación financiera de alta precisión que alcanza un monto ejecutado de \textbf{\$1.199.300~CLP}, lo que representa un \textbf{95,94\% de ejecución presupuestaria}, resguardando un fondo de reserva operacional y contingencias de \textbf{\$50.700~CLP (4,06\%)}. Esta holgura monetaria permite absorber fluctuaciones de flete local, recargos por despacho urgente o sustituciones menores sin desestabilizar la programación financiera.

Para maximizar el impacto ingenieril y asegurar la trazabilidad integral de los desembolsos, el presupuesto se estructura en \textbf{cinco categorías estratégicas}, cada una articulada de manera directa con los Objetivos Específicos (OE) del proyecto:
\begin{enumerate}
    \item \textbf{Categoría 1: Equipos y Componentes Electrónicos de Control (OE1 y OE2):} Financia la unidad microcontrolada central (ESP32), el sistema de censado dual de alta fidelidad (radar de micro-movimiento mmWave de 24~GHz acoplado a detector infrarrojo pasivo PIR), actuadores con aislamiento galvánico y fuentes conmutadas industriales. Cubre la adquisición y diseño de los módulos autónomos para 5 salas piloto más un kit de respaldo (\$279.840~CLP, 23,33\%).
    \item \textbf{Categoría 2: Instrumentación y Medición Energética (OE3):} Provee la infraestructura de metrología eléctrica no invasiva mediante analizadores AC con transformador de corriente (CT) split-core, sistemas de respaldo de registro offline en tarjetas industriales de alta durabilidad, pinza amperimétrica True-RMS certificada por la SEC como patrón de calibración en terreno y dataloggers higrotérmicos autónomos (\$300.580~CLP, 25,06\%).
    \item \textbf{Categoría 3: Materiales de Instalación Eléctrica, Montaje y Seguridad (OE2):} Suministra envolventes estancas IP65 ignífugas para cielo falso, canalizaciones plásticas autoextinguibles, conductores libres de halógenos normalizados bajo pliego técnico SEC RIC, conectores de inserción a resorte WAGO y elementos de protección personal (EPP) dieléctrico de 1000~V (\$293.760~CLP, 24,49\%).
    \item \textbf{Categoría 4: Difusión, Señalética y Documentación Técnica Replicable (OE4):} Financia la transferencia tecnológica y la visibilidad institucional del piloto mediante placas informativas en acrílico cristal grabado láser, manuales impresos de operación y esquemas unilineales para la DGI, pendón roller para divulgación y señalética vinílica para concientización estudiantil (\$180.620~CLP, 15,06\%).
    \item \textbf{Categoría 5: Herramientas de Calibración, Banco de Pruebas y Contingencia (OE1 y OE2):} Asegura la puesta a punto en laboratorio, protecciones selectivas con portafusibles de riel DIN, supresores de transientes de carga inductiva (snubbers RC y varistores) y un maletín de herramientas de crimpado certificado (\$144.500~CLP, 12,05\%).
\end{enumerate}

\subsection{Presupuesto Maestro Consolidado}

La Tabla~\ref{tab:presupuesto_maestro} presenta el Presupuesto Maestro Consolidado de \textit{Campus Zero}. Cada ítem se encuentra debidamente valorizado con precios de mercado vigentes en Chile al año 2026 (IVA incluido), respaldados en cotizaciones formales de proveedores reconocidos del rubro eléctrico y electrónico (MCI Electronics, AFEL, MechatronicStore, Dartel, Gobantes, Sodimac, Imprenta USM y PubliAcrílicos Valparaíso).

\begin{table}[htbp]
\centering
\scriptsize
\renewcommand{\arraystretch}{1.18}
\caption{Presupuesto Maestro Consolidado del Proyecto Campus Zero (Valores en CLP, IVA incluido)}
\label{tab:presupuesto_maestro}
\begin{tabularx}{\textwidth}{>{\centering\arraybackslash}p{0.7cm} >{\raggedright\arraybackslash}p{4.2cm} >{\centering\arraybackslash}p{0.6cm} >{\centering\arraybackslash}p{0.8cm} >{\raggedleft\arraybackslash}p{1.4cm} >{\raggedleft\arraybackslash}p{1.5cm} X}
\toprule
\textbf{ID} & \textbf{Componente / Denominación Técnica} & \textbf{Cant.} & \textbf{Unidad} & \textbf{P. Unit. (\$)} & \textbf{Total (\$)} & \textbf{Justificación Sintética y Proveedor} \\
\midrule
\multicolumn{7}{l}{\cellcolor{gray!20}\textbf{Categoría 1: Equipos y Componentes Electrónicos de Control (OE1, OE2) -- Subtotal: \$279.840 (23,33\%)}} \\
\midrule
1.1 & Microcontrolador ESP32 DevKit V1 (Dual-Core 240MHz, Wi-Fi/BLE) & 6 & un & \$8.990 & \$53.940 & Unidad de procesamiento local autónomo por sala + repuesto lab. MCI Electronics. \\
1.2 & Sensor radar micro-movimiento Doppler mmWave 24~GHz HLK-LD2410B & 6 & un & \$12.990 & \$77.940 & Detección de presencia humana estática y respiración hasta 6m. AFEL Electrónica. \\
1.3 & Sensor PIR gran angular 360$^\circ$/120$^\circ$ HC-SR501 con cúpula Fresnel & 6 & un & \$3.490 & \$20.940 & Detección macro-cinética inicial de ingreso a sala de estudio. MechatronicStore. \\
1.4 & Fuente conmutada compacta AC/DC 220VAC a 5VDC 2A / 10W aislada & 6 & un & \$11.990 & \$71.940 & Alimentación continua y segura de electrónica desde red AC. MCI Electronics. \\
1.5 & Set módulos relé 2CH optoacoplados 10A y shields borneras & 6 & set & \$9.180 & \$55.080 & Conmutación física circuitos iluminación y extractores + interconexión. AFEL/MCI. \\
\midrule
\multicolumn{7}{l}{\cellcolor{gray!20}\textbf{Categoría 2: Instrumentación y Medición Energética (OE3) -- Subtotal: \$300.580 (25,06\%)}} \\
\midrule
2.1 & Analizador AC Peacefair PZEM-004T v3.0 con CT 100A split-core & 6 & un & \$19.990 & \$119.940 & Medición $V$, $I$, $P$ activa (W) y energía (kWh) sin corte de línea. MechatronicStore. \\
2.2 & Kit Datalogger MicroSD SPI + Memoria SanDisk Industrial 32GB & 6 & un & \$11.990 & \$71.940 & Registro de telemetría energética local a 1~Hz independiente de Wi-Fi. MCI. \\
2.3 & Pinza amperimétrica True-RMS UNI-T UT202A+ cert. SEC & 1 & un & \$44.900 & \$44.900 & Instrumento patrón para contraste y calibración de campo. Dartel Electricidad. \\
2.4 & Datalogger higrotérmico USB Elitech RC-4HC con sonda NTC & 2 & un & \$31.900 & \$63.800 & Monitoreo de temperatura y HR para control de ventilación. MechatronicStore. \\
\midrule
\multicolumn{7}{l}{\cellcolor{gray!20}\textbf{Categoría 3: Materiales de Instalación Eléctrica, Montaje y Seguridad (OE2) -- Subtotal: \$293.760 (24,49\%)}} \\
\midrule
3.1 & Gabinete estanco IP65 policarbonato autoextinguible con riel DIN & 6 & un & \$12.990 & \$77.940 & Envolvente ignífuga de seguridad alojada en cielo falso. Dartel Electricidad. \\
3.2 & Canalización plástica 20x10mm y conductor libre halógenos 1,5mm$^2$ & 1 & global & \$112.050 & \$112.050 & 15 canaletas Legrand (30m) y 3 rollos H07Z1-K pliego SEC. Sodimac/Gobantes. \\
3.3 & Accesorios conexión rápida WAGO 221, prensaestopas y anclajes 3M & 1 & global & \$53.870 & \$53.870 & Bornes sin mantención (50 un), 20 prensaestopas IP68 y tarugos/VHB. Dartel/Sodimac. \\
3.4 & Equipamiento Protección Personal (EPP) dieléctrico 1000V y NCV & 2 & set & \$24.950 & \$49.900 & Guantes aislantes clase 00, gafas UV y lápiz detector tensión. Gobantes Electricidad. \\
\midrule
\multicolumn{7}{l}{\cellcolor{gray!20}\textbf{Categoría 4: Difusión, Señalética y Documentación Técnica (OE4) -- Subtotal: \$180.620 (15,06\%)}} \\
\midrule
4.1 & Placas acrílicas cristal grabadas láser y stickers vinílicos & 1 & global & \$99.220 & \$99.220 & 6 placas muro 18x12cm con embellecedores y pack 60 stickers. PubliAcrílicos/USM. \\
4.2 & Pendón Roller retráctil 80x200cm tela PVC mate alta resolución & 1 & un & \$36.900 & \$36.900 & Soporte visual para ferias científicas y presentación ante DGI. PubliAcrílicos. \\
4.3 & Manual técnico de réplica, esquemas unilineales y operación & 5 & un & \$8.900 & \$44.500 & Documento encuadernado couché 170g para transferencia y mantención. Imprenta USM. \\
\midrule
\multicolumn{7}{l}{\cellcolor{gray!20}\textbf{Categoría 5: Herramientas de Calibración, Banco de Pruebas y Contingencia (OE1, OE2) -- Subtotal: \$144.500 (12,05\%)}} \\
\midrule
5.1 & Banco de pruebas prototipado y maletín crimpadora con ferrules & 1 & kit & \$62.800 & \$62.800 & Protoboard, CP2102 USB-UART y crimpadora con 800 terminales. MCI/Sodimac. \\
5.2 & Protecciones eléctricas: portafusibles riel DIN, varistores MOV y snubber & 1 & pack & \$44.800 & \$44.800 & Fusibles 2A/5A, 12 filtros RC y varistores 275VAC supresores inductivos. Gobantes/Mech. \\
5.3 & Kit de repuestos críticos de contingencia inmediata en terreno & 1 & kit & \$36.900 & \$36.900 & Módulo backup completo (ESP32 + mmWave + PIR + Relé) en terreno. AFEL Electrónica. \\
\midrule[\heavyrulewidth]
\multicolumn{4}{r}{\textbf{SUBTOTAL PRESUPUESTO EJECUTADO:}} & \textbf{\$1.199.300} & \textbf{95,94\% de ejecución presupuestaria} \\
\multicolumn{4}{r}{\textbf{FONDO DE RESERVA Y CONTINGENCIAS:}} & \textbf{\$50.700} & \textbf{4,06\% margen de contingencia técnica} \\
\multicolumn{4}{r}{\textbf{PRESUPUESTO TOTAL GENERAL SOLICITADO:}} & \textbf{\$1.250.000} & \textbf{100,00\% tope máximo fondo concursable} \\
\bottomrule
\end{tabularx}
\end{table}

\subsection{Análisis de Viabilidad Financiera, Sostenibilidad y Retorno de la Inversión (ROI)}

El análisis técnico-económico demuestra la sólida viabilidad y rentabilidad del proyecto \textit{Campus Zero}, fundamentado en la eliminación del desperdicio eléctrico pasivo en recintos universitarios:

\subsubsection{Cálculo del Desperdicio Fantasma y Potencial de Ahorro Energético}
Cada una de las salas de estudio seleccionadas en Casa Central cuenta con dos circuitos de consumo desatendido:
\begin{itemize}
    \item \textbf{Circuito 1 (Iluminación):} Cuatro paneles LED de 60$\times$60~cm de 40~W nominales cada uno, totalizando una potencia activa de $0{,}16$~kW.
    \item \textbf{Circuito 2 (Climatización/Ventilación):} Dos extractores forzados de aire de 50~W nominales cada uno, demandando una potencia activa de $0{,}10$~kW.
\end{itemize}
La carga eléctrica simultánea asciende a \textbf{$0{,}26$~kW por sala} a 220~VAC. A partir de las auditorías energéticas preliminares realizadas en el edificio de docencia, se constata un desperdicio fantasma conservador de \textbf{3 horas diarias por sala} en las que luminarias y extractores permanecen encendidos con el recinto completamente desocupado. Esto equivale a:
\begin{equation}
    E_{\text{sala/día}} = 0{,}26\text{ kW} \times 3\text{ h/día} = 0{,}78\text{ kWh/día/sala}
\end{equation}
Considerando un mes calendario estándar de 22 días hábiles lectivos:
\begin{equation}
    E_{\text{sala/mes}} = 0{,}78\text{ kWh/día} \times 22\text{ días} = 17{,}16\text{ kWh/mes/sala} \approx 17{,}2\text{ kWh/mes/sala}
\end{equation}
Para las cinco salas de estudio intervenidas en el piloto, el ahorro mensual totaliza:
\begin{equation}
    E_{\text{piloto/mes}} = 17{,}16\text{ kWh/mes/sala} \times 5\text{ salas} = 85{,}8\text{ kWh/mes}
\end{equation}
En un ciclo académico regular de 9 meses lectivos efectivos (marzo a diciembre, descontando periodos vacacionales):
\begin{equation}
    E_{\text{piloto/año}} = 85{,}8\text{ kWh/mes} \times 9\text{ meses} = 772{,}2\text{ kWh/año académico}
\end{equation}

\subsubsection{Retorno Económico Directo para la USM}
Aplicando la tarifa eléctrica promedio institucional de media tensión suministrada a la Casa Central de la UTFSM en Valparaíso, equivalente a \textbf{\$195~CLP/kWh} (considerando cargos por energía activa y demanda en horas fuera de punta), el retorno financiero directo para el campus en las cinco salas piloto asciende a:
\begin{equation}
    \text{Ahorro Económico Piloto} = 772{,}2\text{ kWh/año} \times \$195\text{ CLP/kWh} = \$150.579\text{ CLP/año}
\end{equation}

\subsubsection{Escalabilidad Institucional a Nivel Campus (50 Salas) y Payback}
El diseño modular del hardware y la entrega de la \textit{Guía Técnica de Replicación y Mantención} (OE4) habilitan el escalamiento progresivo del sistema a la totalidad de las 50 salas de estudio y laboratorios docentes de Casa Central. A esta escala:
\begin{equation}
    E_{\text{campus/año}} = 50\text{ salas} \times 17{,}16\text{ kWh/mes/sala} \times 9\text{ meses} = 7.722\text{ kWh/año}
\end{equation}
\begin{equation}
    \text{Ahorro Económico Campus} = 7.722\text{ kWh/año} \times \$195\text{ CLP/kWh} = \$1.505.790\text{ CLP/año}
\end{equation}
Dado que el costo marginal de hardware por sala se reduce drásticamente al eliminar costos no recurrentes de instrumental y banco de pruebas, el costo de implementación por recinto a escala mayorista desciende a aproximadamente \$85.000--\$95.000~CLP. Esto permite proyectar un \textbf{periodo de recuperación de la inversión (payback simple) de 2,5 a 3 años} para la infraestructura completa de Casa Central, posterior al cual la tecnología genera flujos de ahorro neto recurrente para la universidad.

\subsubsection{Retorno Ambiental y Aporte a la Descarbonización Institucional}
Conforme a los balances oficiales del Coordinador Eléctrico Nacional (CEN) y el Ministerio de Energía, el factor de emisión promedio del Sistema Eléctrico Nacional (SEN) de Chile se sitúa en \textbf{$0{,}35\text{ kg CO}_2\text{e/kWh}$}. Por consiguiente, la mitigación ambiental de \textit{Campus Zero} se cuantifica en:
\begin{itemize}
    \item \textbf{Piloto (5 salas):} $772{,}2\text{ kWh/año} \times 0{,}35\text{ kg CO}_2\text{e/kWh} = \mathbf{270{,}3\text{ kg CO}_2\text{e/año}}$ evitados directamente.
    \item \textbf{Campus General (50 salas):} $7.722\text{ kWh/año} \times 0{,}35\text{ kg CO}_2\text{e/kWh} = \mathbf{2.702{,}7\text{ kg CO}_2\text{e/año}} \approx \mathbf{2{,}7\text{ toneladas de CO}_2\text{e/año}}$ mitigadas en forma permanente, alineándose con las metas de carbono neutralidad institucional de la USM.
\end{itemize}

\subsubsection{Preservación de Vida Útil de Luminarias y Equipamiento Electromecánico}
Al extinguir 3 horas diarias de operación ociosa durante 220 jornadas de actividad universitaria anual (incluyendo periodos de evaluaciones y exámenes), se ahorran más de \textbf{660 horas operativas de vida útil por año} (+660~h/año) en cada panel LED y extractor de ventilación. En luminarias LED, este régimen reduce la degradación térmica del fósforo y el desgaste de los condensadores electrolíticos de los drivers. En los extractores, evita la fatiga prematura de los rodamientos por funcionamiento térmico en vacío y disminuye la acumulación centrífuga de polvo, aplazando las compras de recambio y reduciendo los costos de mantenimiento correctivo de la Dirección General de Infraestructura.

\subsubsection{Sostenibilidad Operacional y Continuidad Post-Proyecto}
Una vez concluido el fondo concursable, el sistema opera con costos de mantenimiento marginal cero. La lógica embebida opera en firmware de código abierto sin licencias recurrentes, la telemetría se almacena localmente y puede enlazarse a servidores internos de la universidad sin requerir plataformas comerciales en la nube, y los repuestos son componentes estándar de alta disponibilidad en el mercado nacional.

\par\noindent\textbf{Articulación con Justificaciones y Cotizaciones:} La justificación técnica exhaustiva, los principios de austeridad, la metodología de cotización comercial y el detalle pormenorizado de los proveedores para cada ítem de esta matriz presupuestaria se presentan en la Sección~\ref{sec:justificaciones}, a cargo del Analista Financiero (Javier). La integración cronológica de las compras e instalaciones con las 11 actividades del proyecto se expone en la Sección~\ref{sec:gantt} (Jose).